\section{Related Work}\label{sec:related}
\textbf{Cryptographic capability of agents.} ACTS~\cite{acts} introduced a
multi-tier classification of LLM cipher identification under controlled blind
conditions and reported a 40.9-point metadata-aided vs.\ blind gap and a
92.9\% vs.\ 41.3\% agentic vs.\ tool-only contrast. AICrypto~\cite{aicrypto}
evaluates cryptographic knowledge, CTF exploitation, and proofs, while
CryptanalysisBench~\cite{cryptanalysisbench} tests attacks on cryptographic
primitives. These are direct capability comparators, not merely classification
benchmarks. FirmAgent~\cite{firmagent} and APPATCH~\cite{appatch} show that LLM-based
workflows can discover or patch real vulnerabilities, but their objectives and
oracles differ from cryptographic attack construction. CyberGym~\cite{cybergym} evaluates executable vulnerability
reproduction against real codebases; AutoPT~\cite{autopt} uses state-machine
constraints for web penetration testing. AutoPenBench~\cite{autopenbench}
records intermediate milestones and distinguishes autonomous from human-assisted
testing. Our distinguishing measurement holds recovery instances and execution
conditions fixed for representation contrasts, then holds generated bytes
fixed for a separate evaluator-policy comparison; the object is the relation
between construction, acceptance and attribution, not broader attack coverage.

\textbf{Attack mechanisms.} Cache-timing attacks on
AES~\cite{bernstein,osvik,liu15,gruss16,gullasch11} and remote timing
attacks~\cite{brumley03,kocher} establish the side-channel primitives we probe.
Protocol-failure attacks---FREAK~\cite{freak}, Logjam~\cite{logjam},
POODLE~\cite{poodle}---and hybrid key-exchange design and
deployment~\cite{bindel,paquin2020,rfc9954,rfc10024,nist2024} define the
protocol and deployment weaknesses in our target set. We introduce no new
attacks; we measure whether agents can execute these known classes.

\textbf{Trajectories and evaluation integrity.} CTF-ABACUS~\cite{ctfabacus}
provides trace-level provenance for agentic offensive-security evaluation;
ABC~\cite{abc} audits task setup and reward design, while
AgentSuite~\cite{agentsuite} distinguishes user, environment, ground-truth,
and evaluation defects. BenchJack~\cite{benchjack}, a preprint, audits benchmark
hackability. AgentRewardBench~\cite{agentrewardbench} compares automatic
trajectory evaluations with expert judgments; PROV-AGENT~\cite{provagent}
tracks workflow provenance. ABC demonstrates score distortion from task and reward defects, including
rescoring fixed agent artifacts; AgentSuite diagnoses validity by benchmark
component. ACTS-Attack adds a controlled cryptographic mechanism analysis:
the recovery instance is matched across representations, generated programs
are fixed across admission rules, and output-preserving cleanup tests the
dominant source of the measured distortion. AgentDojo~\cite{agentdojo}, IsolateGPT~\cite{isolategpt},
and SAGA~\cite{saga} address agent execution isolation and security
architecture; ToolHijacker~\cite{toolhijacker} demonstrates that tool-selection
interfaces are themselves attack surfaces. ORTHRUS~\cite{orthrus} illustrates the
separate provenance problem of making attack evidence attributable and reviewable.
Cloak, Honey, Trap~\cite{cheat} shows that target-side deception can disrupt
offensive agents, making the agent's own execution environment part of the
security boundary. The attribution experiments complement this mechanism analysis by testing
which authenticated, mutually consistent records support a cryptographic
execution claim. Existing provenance mechanisms can implement the policy;
the contribution is its operational specification and tested behavior under
the stated trust model.

\textbf{Safety evaluation.} AgentHarm~\cite{agentharm} evaluates refusal and harmfulness. Our
safety probe (SCT, B1 ladder) measures whether alignment changes probe-execution
behavior on cryptographic tasks, not general harmfulness.

\textbf{Comparison instrument.} We audit six agent-security releases---Cybench, NYU CTF Bench, CVE-Bench, Incalmo, AgentHarm, and WebSP-Eval~\cite{cybench,nyuctf,cvebench,incalmo,agentharm,webspeval}---with a predeclared 9-field rubric (Section~\ref{sec:audit}). This supplies release-level context for, but does not replace, the attack evidence.
